% MG01. Inserción: después de exc:aw y su interpretación elíptica;
% antes de exc:codigo. Requiere L_1 y la carta A_W ya construidos.
\subsubsection{Neutralité duale et orientation de la première frontière conjointe}
\label{mg01:frontera}

La frontière de prolongement réunit les trois biographies régionales par
une relation qui conserve l'axe, l'agrégat latéral, la saturation et
l'occupation des colonnes. Dans la coordinatisation par vecteurs-lignes déjà fixée,
les derniers blocs de longueur six des mots de longueur trente
sont
\[
 u_4^\pi=110221,\qquad u_4^e=102011,\qquad
 u_4^\varphi=010200.
\]
Les exposants conservent les trois régions de
\eqref{eq:palabras-regionales} ; chaque mot s'interprète comme un
élément de \(\mathbb F_3^6\), avec des représentants entiers dans
\(\{0,1,2\}\). La relation de frontière transporte l'axe au moyen de
\begin{equation}
 z:=u_5^\varphi=u_4^e=102011,
 \qquad a:=u_4^\varphi A_WL_1^3=111101.
 \label{mg01:eje-contraste}
\end{equation}
Ici, \(a\) est le contraste orienté entre les deux canaux latéraux et
l'axe. L'agrégat latéral correspondant est
\begin{equation}
 x+y=210112,\qquad
 q:=x+y+z=012120,\qquad a=x+y-z=111101
 \quad\text{en }\mathbb F_3^6,
 \label{mg01:agregados}
\end{equation}
où \(x=u_5^\pi\) et \(y=u_5^e\). La somme et le contraste ont des
lectures duales distinctes : \(qA_W=012000\) et \(aA_W=120210\).

Soit \(B\in M_{3,6}(\mathbb F_3)\) la matrice de lignes \(x,y,z\).
Son relèvement canonique permet de définir, colonne par colonne,
\[
 c_j=\left\lfloor\frac{x_j+y_j+z_j}{3}\right\rfloor,
 \qquad
 o_j=\mathbf1_{x_j\ne0}+\mathbf1_{y_j\ne0}+\mathbf1_{z_j\ne0}.
\]
% BEGIN R28_BOUNDARY_OCCUPANCY_INPUT
% BEGIN R28_BOUNDARY_OCCUPANCY
% Insertar en MG01 tras la definición de c_j,o_j, antes de su valor.
La loi de la frontière de profondeur six fixe \(c_j=1\). L'axe \(z\) et l'
agrégat \(q\) étant déjà déterminés, les sommes entières latérales sont
\[
 (x_j+y_j)_{j=1}^6=(2,4,3,4,4,2).
\]
La règle d'occupation minimale sélectionne, parmi ses décompositions avec \(x_j,y_j\in\{0,1,2\}\), les colonnes ayant le plus petit nombre d'entrées non nulles, en comptant également \(z_j\). Pour une somme de deux, les paires extrêmes \((0,2),(2,0)\) occupent une position latérale ; pour une somme de trois, les paires \((1,2),(2,1)\) en occupent deux ; pour une somme de quatre, seule existe \((2,2)\). En ajoutant l'occupation de l'axe, on obtient \(o=(2,2,3,2,3,2)\). À cette frontière, avec \(s=[x+y]_3\), la même règle s'exprime comme
\[
 o_j=2+\mathbf1_{\{z_j\ne0,\ s_j\ne2\}}.
\]
% END R28_BOUNDARY_OCCUPANCY

% END R28_BOUNDARY_OCCUPANCY_INPUT

À cette frontière, la saturation et l'occupation sont
\begin{equation}
 c=111111,\qquad o=223232.
 \label{mg01:saturacion}
\end{equation}
L'égalité entière de chaque colonne est donc \(x_j+y_j+z_j=q_j+3\). La retenue \(c\) enregistre le relèvement entier de cette frontière ; son domaine est celui des six colonnes présentes.

\begin{proposition}[Huit distributions d'une frontière saturée]
\label{mg01:ocho}
Les conditions \eqref{mg01:eje-contraste}--\eqref{mg01:saturacion}
laissent exactement huit matrices. Leurs alternatives par colonnes sont
\begin{equation}
\begin{array}{c|c|c|c}
 j&z_j&x_j+y_j\ \text{en }\mathbb Z&(x_j,y_j)\\\hline
 1&1&2&(0,2),(2,0)\\
 2&0&4&(2,2)\\
 3&2&3&(1,2),(2,1)\\
 4&0&4&(2,2)\\
 5&1&4&(2,2)\\
 6&1&2&(0,2),(2,0)
\end{array}
\label{mg01:tabla-columnas}
\end{equation}
\end{proposition}
\begin{proof}
La troisième ligne est fixée par le transport de l'axe. L'égalité
entière \(x_j+y_j=q_j+3-z_j\), jointe à \(o_j\), donne les alternatives
du tableau. Les première, troisième et sixième colonnes admettent deux choix
indépendants ; les trois autres en admettent un. Le produit de leurs cardinaux
est \(2\cdot1\cdot2\cdot1\cdot1\cdot2=8\), et chaque choix conserve
toutes les égalités déclarées.
\end{proof}

La neutralité est évaluée dans la coordinatisation duale. Avec
\(\mathbf1=(1,1,1,1,1,1)^{\mathsf T}\), la matrice de
\eqref{exc:aw} détermine
\begin{equation}
 \omega=A_W\mathbf1=(2,1,1,1,1,1)^{\mathsf T},
 \qquad A_W\omega=-\mathbf1.
 \label{mg01:unidad-dual}
\end{equation}
La condition de clôture duale de la frontière est
\begin{equation}
 (BA_W)\mathbf1=B\omega=0\quad\text{en }\mathbb F_3^3.
 \label{mg01:neutralidad}
\end{equation}
Cette condition agit sur la collection des frontières saturées. La relation
quadratique de \(A_W\) détermine le transport de l'unité duale ; la
neutralité spécifie quelles frontières satisfont la clôture correspondante.

\begin{proposition}[Paire duale et frontière orientée]
\label{mg01:dos-orientacion}
La neutralité \eqref{mg01:neutralidad} laisse exactement
\begin{equation}
 B_-=
 \begin{pmatrix}0&2&1&2&2&2\\2&2&2&2&2&0\\1&0&2&0&1&1\end{pmatrix},
 \qquad
 B_+=
 \begin{pmatrix}2&2&2&2&2&0\\0&2&1&2&2&2\\1&0&2&0&1&1\end{pmatrix}.
 \label{mg01:pareja}
\end{equation}
L'involution qui échange clôture et propagation permute \(B_-\)
et \(B_+\). Dans la coordinatisation semi-ouverte avec l'orientation APP fixée pour
les canaux, la frontière canonique est \(R_{36}=B_+\).
\end{proposition}
\begin{proof}
Nous écrivons les trois choix binaires de la première ligne sous la forme
\[
 x=(2\varepsilon_1,2,1+\varepsilon_3,2,2,2\varepsilon_6),
 \qquad \varepsilon_1,\varepsilon_3,\varepsilon_6\in\{0,1\}.
\]
La condition \(x\omega=0\) équivaut à
\[
 1+\varepsilon_1+\varepsilon_3-\varepsilon_6\equiv0\pmod3.
\]
Les possibilités sont \((0,0,1)\) et \((1,1,0)\). La ligne de l'axe
satisfait \(z\omega=0\) ; l'agrégat latéral a également une évaluation
duale nulle. Par conséquent, dans les deux cas, \(y\omega=0\). On obtient
exactement les deux matrices de \eqref{mg01:pareja}.

Le choix final applique l'orientation APP de la coordinatisation semi-ouverte à
cette paire spéculaire. Dans cette coordinatisation, les représentations ordonnées sont
\[
 B_-\longmapsto(789,048,894),\qquad
 B_+\longmapsto(793,045,894).
\]
Le représentant de l'orientation fixée est \(B_+\). Cette opération de sélection conserve l'ordre des trois régions et la convention de frontière de leurs cylindres. La réduction algébrique de huit matrices à deux et le choix du représentant orienté constituent ainsi deux étapes de la même construction.
\end{proof}

Les charges visibles distinguent ces étapes avec précision :
\begin{equation}
 B_-\mathbf1=(0,1,2)^{\mathsf T},\qquad
 B_+\mathbf1=(1,0,2)^{\mathsf T},\qquad
 B_\pm\omega=0.
 \label{mg01:cargas}
\end{equation}
La charge visible et la neutralité duale sont des évaluations sur des vecteurs
distincts. La frontière sélectionnée conserve trois lignes de six trits,
d'indices ambiants \((726,215,301)\), et fournit l'état
de raccordement \(R_{36}\) du prolongement conjoint. Sa mémoire inclut
l'axe transporté, la saturation, l'incidence des colonnes et
l'orientation qui distingue les feuillets de la paire.
